\section{Constructing the harmonic lift}
\label{sec:harmonic-construction}

We now start in characteristic $23$, rather than reduce a known generic
equation.  The purpose is to produce a cover with the osculating
incidence.  The construction has three stages: lift a special $M_{23}$-map,
derive its marked canonical neighbourhood, and identify the lift by
exact local uniqueness.  The first two stages are developed here; the
last is the local reconstruction theorem below.

Throughout this section put
\[
 k=\overline{\FF}_{23},\qquad K^{\mathrm{ur}}=\operatorname{Frac}(W(k)),
 \qquad K=K^{\mathrm{ur}}(\pi),\qquad \pi^2=-23,
\]
and let $\mathcal O$ be the valuation ring of $K$, with $v(\pi)=1$.
Write $G=M_{23}$, let $M=M_{22}$ be a point stabilizer, and choose
$P\simeq C_{23}$ and $D=N_G(P)=P{:}H$, where $H=M\cap D\simeq C_{11}$.
The letter $D$ in this section denotes this normalizer, not a residual
point in the divisor relation of \cref{thm:osculating-summary}.

\subsection{The special map and its lift}

A \emph{tail} is a component meeting the rest of a semistable target
in one point.  A primitive tail contains a specialization of an
original tame branch point; a new tail contains none.  The associated
covers below are separable degree-$23$ maps; their normal closures,
not the quotient maps themselves, are $G$-covers.

\begin{computation}[Two characteristic-$23$ tail covers]
\label{comp:harmonic-tail-groups}
Over $\FF_{23}$ consider
\begin{equation}
 \begin{aligned}
  D_7:&\quad T=X^{23}+7X^9+18X^2,\\
  D_{15}:&\quad w^3=19(X^5+17),\qquad T=(X^6+2X)w.
 \end{aligned}
\label{eq:harmonic-tail-inputs}
\end{equation}
Their normal closures have geometric Galois group $M_{23}$.
At infinity the inertia groups are $23{:}11$, with lower jumps $7$
and $15$, respectively.  The first source has genus zero and finite
branch fibre $2^8 1^7$; the second has genus four and is \'etale over
the affine target.
\end{computation}

Here the group assertion is a function-field computation, recorded
separately for the two equations in \cref{sec:certificates}.
The arithmetic groups over $\FF_{23}(T)$ are $M_{23}$.  Their geometric
subgroups are normal and contain nontrivial inertia, so simplicity
gives the same geometric groups.  The remaining assertions can be
checked directly.  For the first map,
\[
 T=(X^8+17X)^2(X^7+12),\qquad T'=17(X^8+17X),
\]
with squarefree, coprime factors.  Its different at infinity has
degree $44-8=36$.  For the second, put $H_1=X^6+2X$ and $Q_1=X^5+17$.
The identity $3H_1'Q_1+H_1Q_1'=10$ proves \'etaleness away from
$w=0$; there $w$ is a parameter and $H_1$ a unit.
The pole orders of $X,w,T$ are $3,5,23$, and the different has degree
$52$.  The wild inertia is $23{:}m$ with $m\mid11$.
If $m=1$, the quotient different would be divisible by $22$.
Thus $m=11$, and the formula $22+2h$ for that different gives the
stated jumps.

On $Z_0=\PP^1_a$ set
\begin{equation}
 \omega=-\frac{a^6}{a^{22}-1}\,da,
 \qquad x=\frac{2a^{11}}{a^{11}-1}.
\label{eq:harmonic-deformation-datum}
\end{equation}
The residue at $c\in\FF_{23}^{\times}$ is $c^7$.  Consequently
$\omega=d\log\prod_c(a-c)^{[c^7]}$, where the brackets choose
integer representatives between $1$ and $22$.
If $\chi:H\to\FF_{23}^{\times}$ is conjugation on $P$, let $H$ act
on $a$ through $\chi^8$.  Since $7\cdot8\equiv1\pmod{11}$,
$h^*\omega=\chi(h)\omega$.
The ramification data $(m,h)$ at $x=0,1,\infty,2$ are
\[
 (11,7),\quad(1,0),\quad(1,0),\quad(11,15).
\]
There are exactly three ratios $h/m<1$, none equal to $1$, and all
are less than $2$.

Choose a wild point of each normal closure, identify its inertia with
$D$, and its tame complement with $H$.  Attach these two full $G$-tails
to the induced central curve $\operatorname{Ind}_D^G(Y_0)$, where
$Y_0\to Z_0$ is purely inseparable of degree $23$.
The matched inertia groups and conductors, together with generation
of $G$ by the connected tail groups, verify the special
$G$-deformation datum conditions of \cite[Definition~2.13]{Wewers}.
This constructs a connected special $G$-map, to which
\cite[Theorem~4.5]{Wewers} applies.

It is useful to retain the inseparable map explicitly.  If $a=b^{23}$
on $Y_0$ and $z=b^{11}$ on $N=Y_0/H$, then
\begin{equation}
 x=\frac{2z^{23}}{z^{23}-1},\qquad
 \overline\beta=\frac{1}{1-x}
              =\frac{1-z^{23}}{1+z^{23}}.
\label{eq:harmonic-central-map}
\end{equation}
The primitive attachment, new attachment, and two wild marks are
$e=0,q=\infty,A=1,B=-1$.  Their harmonic position is therefore part
of the constructed special map.  Replacing $a=b^{23}$ by $a=b$ would
give the wrong quotient map.

The residues in the two pole orbits $a^{11}=1,-1$ have opposite
quadratic characters.  They give the opposite classes $23A,23B$ in
the lift.  A common change of wild generator rescales both classes
together; taking inverses preserves square class.  The tame tail
supplies the class $2A$.

\begin{proposition}[Descent of the marked lift]
\label{prop:harmonic-marked-descent}
After ordering the two wild branches, the constructed special map
has one isomorphism class of lifts over $\overline K$, and this class
has a $G$-cover representative over $K$.  Its degree-$23$ quotient
has two individually $K$-rational totally ramified points $A,B$.
\end{proposition}

\begin{proof}
We compute the automorphisms of this particular special map directly.
On $D_7$, a pointed source automorphism is $X\mapsto aX+b$.
It preserves the eight critical points, whose sum is zero, so $b=0$.
Comparison of coefficients gives $a^{14}=a^{21}=1$, or $a^7=1$.
All seven scalings preserve the map up to target scaling.

On the smooth completion of $Y^3=X^5+\lambda$, $\lambda\ne0$, let
$q$ be infinity.  Since $L(3q)=\langle1,X\rangle$ and
$L(5q)=\langle1,X,Y\rangle$, a pointed automorphism has the form
\[
 X\longmapsto aX+b,\qquad Y\longmapsto cY+dX+e.
\]
The $Y^2$ and $X^4$ coefficients force $d=e=b=0$, and then
$a^5=c^3=1$.  These fifteen automorphisms preserve
$Y(X^6+15\lambda X)$ up to scaling.  Their action on the parameter
$Y/X^2$ at $q$ is faithful.

The group facts $\Out(G)=1$, $N_G(M)=M$, and $N_G(D)=D$ identify
these quotient symmetries with the pointed $G$-equivariant tail
symmetries.  Indeed, a quotient symmetry lifts to the normal closure;
correct its action on $G$ by an inner automorphism.  The correcting
element normalizes $M$, so does not change the quotient symmetry.
The point with stabilizer $D$ is unique in its orbit, hence is fixed.

A $G$-equivariant automorphism of the special map preserves the chosen
central component, and its restriction there over the original target
lies in $H$.  Such an element fixes both attaching points.  Each
pointed tail symmetry fixes all their $G$-translates.  Thus these
choices glue independently, and restriction to the normalization
recovers them.  The special-map automorphism group has order
\[
 |H|\cdot7\cdot15=1155.
\]
No matching of derivatives is required to glue automorphisms of the
nodal special fibre $xy=0$.  This argument does not use a general
automorphism-count formula for special $G$-maps.

By \cite[Proposition~4.4 and Theorem~4.5]{Wewers}, the lifts with
identified special fibre form one tame inertia orbit of size
$22\cdot7\cdot15=2310$.  Forgetting that identification divides by
the group just computed, freely: a stabilizer would give a generic
$G$-equivariant automorphism over the target, belonging to $Z(G)=1$.
There are two unframed classes, exchanged by the quadratic cyclotomic
action.  Fixing the ordering of $23A,23B$ leaves one class over $K$.
The $G$-equivariant isomorphisms with its conjugates are unique, so
satisfy the descent cocycle condition.  They descend the cover, not
merely its moduli point.  Taking the $M$-quotient gives the source
over $K$.  Each totally ramified fibre contains one geometric point,
which is consequently $K$-rational.
\end{proof}

\subsection{The marked distance determines the tame character}

The next lemma explains why the abstract good-reduction curve is not
the only information needed.  Distances are logarithmic annulus
lengths, measured with $v(\pi)=1$.  Equivalently, in a smooth residue
disk the distance to the disk separating two sections is the
valuation of their parameter difference.

\begin{lemma}[Marked distance and tame descent]
\label{lem:marked-distance}
Let $K'$ be a complete discretely valued field with algebraically
closed residue field of characteristic $p$.  Let $C/K'$ be a smooth
proper curve of genus at least two with potentially good reduction
over a tame extension.  Suppose that distinct $K'$-points $A,B$
reduce to $q$ on the good geometric special fibre $D'$, and that
$\Aut(D',q)$ is cyclic of order $n>1$, prime to $p$, with faithful
tangent character.  If the marked separating vertex has distance
$1/n$ from the good-reduction vertex, then the minimal good-reduction
extension has degree $n$.  Writing it as $K'(\rho)$ with
$\rho^n=\pi$, the action $\tau(\rho)=\zeta\rho$ has tangent
character $\zeta$ at $q$.
\end{lemma}

\begin{proof}
Over a finite tame Galois extension with good reduction, the descent
action extends by uniqueness of the stable model.  Its kernel
$\Delta$ on $D'$ can be removed without losing smoothness.
At any special point the completed ring is $R'[[u]]$.
Averaging $u$ over $\Delta$ gives an invariant parameter $u_0$
with the same reduction.  Thus the invariant ring is
$R'^{\Delta}[[u_0]]$, by coefficientwise invariance.  The quotient
of the proper projective model exists and is smooth by these local
descriptions.  Minimality therefore makes the residual action
faithful; its order $m$ divides $n$.

For an integral parameter $\xi$ at $q$, the marked distance says
$v(\xi(A)-\xi(B))=1/n$.  Both values lie in the degree-$m$
good-reduction field, whose value group is $\frac1m\ZZ$.
Hence $n\mid m$, so $m=n$.
This difference valuation is unchanged by an integral parameter
change: the quotient of the two differences is its unit linear
coefficient plus terms in the maximal ideal.

Tameness and the algebraically closed residue field give
$K'(\rho)$ as stated.  First linearize a special parameter and then
lift it using the character projector
$n^{-1}\sum_j\zeta^{-ej}\tau^j$.  This projector is linear over
the base valuation ring, not over the extension ring; its reduction
is still a parameter.  We obtain $\tau(\xi)=\zeta^e\xi$ exactly.
Since the sections are rational, their values belong to
$\rho^e\mathcal O_{K'}$, for $0\le e<n$.  The nonzero difference
has valuation $e/n$ plus an integer.  Equality to $1/n$ forces $e=1$.
\end{proof}

In our application $23\nmid|M|$, so the quotient of the semistable
$G$-model by $M$ commutes with reduction.  Since $D$ is transitive
on the $23$ letters, the quotient has one central component and one
node at each attachment.  Its three components form a chain
\[
 D_7\;\text{(rational)}\quad--\quad
 N\;\text{(marked rational)}\quad--\quad D_{15}\;\text{(genus four)}.
\]
Forgetting the map and markings contracts the rational tree and
gives good reduction $D_{15}$.
The target annulus depths are $23/(2h)$ in $v_{23}$ units.
Each source annulus maps with degree $23$; its length is the target
length divided by $23$ \cite[Corollary~2.6]{BPR}.
In our valuation the two source lengths are therefore $1/7$ and
$1/15$.  Dividing by $253$, the order of the inertia in the Galois
closure, would be incorrect for this quotient map.

The pointed automorphism calculation above and
\cref{lem:marked-distance} give the good-reduction field
$K(\rho)$, $\rho^{15}=\pi$.  On $Y^3=X^5+\lambda$ its descent
action is
\begin{equation}
 \tau(X)=\zeta^{-3}X,\qquad \tau(Y)=\zeta^{-5}Y.
\label{eq:harmonic-good-character}
\end{equation}
Indeed this is the unique pointed automorphism whose character on
$Y/X^2$ is $\zeta$.  The hypotheses of the lemma are essential:
neither an unmarked good-reduction curve nor a wild extension would
justify this argument.

\subsection{A canonical model from differentials}

\begin{lemma}[The descended canonical lattice]
\label{lem:harmonic-canonical-lattice}
The marked quotient of \cref{prop:harmonic-marked-descent} has a flat
canonical complete-intersection model over $\mathcal O$ with regular
total space.  Its closed fibre is rational with one cusp of value
semigroup $\langle3,5\rangle$.  In suitable closed coordinates it is
\[
 HW-U^2-\alpha UV-\beta V^2=0,\qquad UW^2-V^3=0.
\]
Its affine normalization is
$(H,U,V)=(z^6+\alpha z^4+\beta z^2,z^3,z)$.
\end{lemma}

\begin{proof}
On the good special fibre the canonical forms
\[
 h=\frac{dX}{Y^2},\quad u=Xh,\quad v=Yh,\quad w=X^2h
\]
have vanishing orders $6,3,1,0$ at infinity and characters $7,4,2,1$
modulo $15$.  The relative dualizing sheaf is invertible and commutes
with base change \cite[Tag~0E6N]{Stacks}.
Lift these forms by cohomology and base change, and
apply the tame character projectors.  Their reductions form a basis,
so their lifts do too.  Multiplication in degrees two and three gives
equivariant canonical equations $q_g,c_g$ of characters $8,6$ reducing to
\[
 hw-u^2=0,\qquad uw^2-v^3+\lambda h^3=0.
\]
Surjectivity of these multiplication maps and the canonical Hilbert
function identify the relative ideal with this quadric and cubic.
Only canonical forms have been lifted.  Lifting $X,Y$ themselves
with their special pole orders would impose extra Weierstrass
conditions on the generic curve.

Set $(H,U,V,W)=(h,\rho^3u,\rho^5v,\rho^6w)$.
All four coordinates have character $7$; dividing them all by
$\rho^7$ gives an invariant generic basis.  Thus the embedding is
into $\PP^3_K$, with no unexamined twist of the projective space.
Its equations are
\begin{equation}
 \begin{aligned}
 q&=\rho^6q_g(H,\rho^{-3}U,\rho^{-5}V,\rho^{-6}W),\\
 c&=\rho^{15}c_g(H,\rho^{-3}U,\rho^{-5}V,\rho^{-6}W).
 \end{aligned}
\label{eq:harmonic-canonical-lattice}
\end{equation}
For a monomial of weight $r$, using weights $(0,3,5,6)$, the
good-model coefficient has $\rho$-order congruent to $r-6$ for $q$
and $r-15$ for $c$.  If $n_q,n_c$ are their least nonnegative
residues, then $n_q+6-r$ and $n_c+15-r$ are nonnegative for all ten
quadratic and twenty cubic monomials.  This proves integrality of
every coefficient.  The possible closed equations are precisely
\[
 \begin{aligned}
 HW-U^2+aUV+bUW+cV^2+dVW+eW^2&=0,\\
 UW^2-V^3+fV^2W+gVW^2+hW^3&=0.
 \end{aligned}
\]
On $W=1$ put $V=z$ and solve for the monic cubic $U$ and monic
sextic $H$.  Translate $z$ to remove the quadratic term of $U$;
linear changes remove its linear and constant terms and the cubic,
linear and constant terms of $H$.  The degree-five term of $H$
vanishes with the same translation.  This gives the asserted form.

The affine chart is $\mathbb A^1$.  At infinity the coordinate orders
are $3,5,6$, and the arithmetic genus of the complete intersection is
four.  Thus the unique singularity has $\delta=4$ and semigroup
$\langle3,5\rangle$, whose gaps are $1,2,4,7$.
The closed equations form a regular sequence, giving flatness of
the relative intersection; adjunction gives $\omega=\mathcal O(1)$.
Finally the $H^3$ coefficient of $c$ is $\pi$ times a unit reducing
to $\lambda$.  At the cusp, on $H=1$, the linear parts of the two
equations in the total local ring are $W+a_0\pi$ and $\lambda\pi$.
They are independent.  The total surface is regular there, and is
relatively smooth at every other closed point.
\end{proof}

In the good model's disk at $q$, the lifted ratio $\xi=v/w$ is an
integral parameter of character one.  The weighted coordinates give
$z=\xi/\rho=V/W$; their other affine ratios are integral power
series on the disk $|\xi|\le|\rho|$ and reduce to the cubic and sextic
above.  This is the marked central disk: $\xi(A)\in\rho\mathcal O$,
and the separating radius is $|\rho|$ by
\cref{lem:marked-distance}.  Thus its reduction is $N$ of
\cref{eq:harmonic-central-map}, not an unrelated rational curve.

\begin{proposition}[The effective divisor fixes the closed equations]
\label{prop:harmonic-effective-differential}
Choose the harmonic coordinate on $N$ with $e=0,q=\infty,A=1,B=-1$.
The closure $R_8$ of the eight simple ramification points is finite
flat of degree eight, with closed divisor $8e$ on the canonical
model.  The divided logarithmic differential is integral and satisfies
\begin{equation}
 \overline\eta=\overline{\frac{d\log\beta}{23}}
       =\frac{2z^8}{z^2-1}\,dz.
\label{eq:harmonic-effective-differential}
\end{equation}
The closed equations can consequently be taken to be
$q_0=HW-U^2-UV-V^2$ and $c_0=UW^2-V^3$ in this marked coordinate.
\end{proposition}

\begin{proof}
After a finite extension splitting the ramification points, all eight
lie in the primitive disk and have $z$ of positive valuation.
The integral canonical ratios just constructed reduce at those
points to their values at $e$.  Properness and the valuative criterion
therefore locate all their centers on the canonical model at $e$.
This argument does not assume a contraction morphism from the
semistable model to the canonical model.

Their schematic closure is proper with finite fibres, hence finite.
It is torsion-free over $\mathcal O$, hence flat of degree eight.
Near $e$ the surface is smooth, so this is an effective Cartier
divisor with closed equation $z^8$.  Retaining this multiplicity is
essential in the next step.

On the regular surface, $\operatorname{div}(\beta)=23A-23B$:
there is no vertical contribution because the central reduction is
nonconstant.  Normality makes $\beta$ a unit away from these disjoint
sections and extends it to a morphism to $\PP^1$.  In particular
there is no hidden base point at the cusp.

Take a primitive integral multiple $s$ of $d\log\beta$ in
$\omega(A+B)$.  Its restriction to $R_8$ is generically zero and
therefore zero, since the latter scheme is flat.  Its closed-fibre
restriction is nonzero: otherwise the multiplication-by-$\pi$ exact
sequence would contradict primitivity.  On the normalization,
\[
 \nu^*\omega(A+B)=\omega_{\PP^1}(8q+A+B)
\]
has degree eight.  Its section already vanishes at $8e$, so it has
no further zero, and has nonzero residues at both marks.  The
residues of $d\log\beta$ are $23,-23$.  Thus the primitive scaling
divides by exactly $23$, up to a unit.  Normalizing the residue at
$A$ to one gives \cref{eq:harmonic-effective-differential}.

It remains to normalize the canonical equations without moving the
four directions.  Subtract $d\log((z-1)/(z+1))$ from the displayed
differential.  This leaves the dualizing differential
\[
 2(z^6+z^4+z^2+1)\,dz.
\]
The subtracted differential is regular at $q$ on the normalization,
so belongs to the dualizing sheaf there.  The canonical section $W$
has no affine zero and is a constant multiple of $dz$.
Using the canonical degrees $0,1,3,6$, write the two monic higher
polynomials as $z^3+u_0z^2$ and
$z^6+v_5z^5+v_4z^4+v_2z^2$ after removing lower basis terms.
The quadric gives $v_5=2u_0$; the displayed differential then gives
$u_0=v_5=0$ and $v_4=v_2=1$.  This proves the stated equations in
the harmonic coordinate, without simultaneously postulating an
unjustified polynomial and marked-point normalization.
\end{proof}

\subsection{Comparing the two tails on the same sheets}

The first necessary deformation equations have two solutions.
Only one is compatible with the common $G$-action.  To compare them,
one must transport the same wild generator through both nodes.
The following elementary calculation makes that transport precise.

\begin{lemma}[A leading coefficient at a node]
\label{lem:harmonic-node-coefficient}
Let $S$ be a complete discrete valuation ring with uniformizer $t$,
let $\alpha\in S$ have valuation $e>0$, and put
$\mathcal A=S[[a,y]]/(ay-\alpha)$.
Write $v_C,v_D$ for the vertical valuations of the branches $y=0$
and $a=0$, normalized by $v_C(t)=v_D(t)=1$.
Suppose $h\ge2$, $F\in\mathcal A[1/t]$, and
\[
 v_C(F)\ge he,\qquad v_D(F)\ge e,\qquad
 (F/\alpha^h)|_C=c a^{1-h}+O(a^{2-h}).
\]
Then $(F/\alpha)|_D=c y^{h-1}+O(y^h)$.
In particular, if $\sigma$ preserves the node and its branches and
$F=\sigma(a)-a$ satisfies the central condition, then
\[
 \sigma(y)-y=-c y^{h+1}+O(y^{h+2}).
\]
\end{lemma}

\begin{proof}
The normal form in $a,y$, with a bounded $t$-denominator, gives a
Laurent expansion $F=\sum f_na^n$ after substituting $y=\alpha/a$.
The two valuation bounds say coefficientwise
$v(f_n)\ge he$ and $v(f_n)+en\ge e$.
There is no cancellation between distinct powers at either branch;
only finitely many negative powers occur at a fixed $t$-order.
The $n$th contribution to $F/\alpha$ on $D$ is
$f_n\alpha^{n-1}y^{-n}$.  It vanishes for $n>1-h$ and has
coefficient $c$ for $n=1-h$.  Lower indices give higher powers.
For an automorphism, the second valuation bound follows from
$v_D(\sigma(a))=v_D(a)=e$, and
\[
 \sigma(y)-y=-(F/\alpha)y\sigma(y)
\]
gives the final assertion.
\end{proof}

We verify that the coordinates used here are actual node parameters.
The primitive quotient disk has $z$-radius $|\pi|^{1/7}$.
It is Galois invariant and contains a point over a tame extension:
lift a smooth point of its component.  Averaging the conjugates of
its coordinate gives a $K$-rational center in the same disk.
That center reduces to $0$, hence lies in $\pi\mathcal O$ and is
strictly inside the radius in question.  The disk is therefore
centered at $0$.  The primitive open annulus is
$|\pi|^{1/7}<|z|<1$.  The new open annulus is
$|\rho|<|\xi|<1$, with coordinate $1/z=\rho/\xi$.
Neither coordinate has a horizontal zero or pole in its annulus.

On a chosen $G$-node the quotient to the degree-$23$ source is by
$H=C_{11}$.  Tame averaging gives an eigenparameter $x$, with the
opposite parameter of inverse character.  The invariant node has
parameters $x^{11},y^{11}$.  The pullback of $z$, or of $1/z$, has
the same vertical divisor as $x^{11}$ and no horizontal divisor
through the node.  Their ratio is a unit in the normal completed
ring; an eleventh root of that unit exists by Hensel's lemma.
We may consequently choose actual parameters $a_7,a_{15}$ with
\[
 a_7^{11}=z,\qquad a_{15}^{11}=1/z.
\]
These coordinates persist in the formal projective slice used below:
a change $I+\pi M_0$, conjugated by the weights $(0,3,5,6)$, changes
the good coordinates only in $\rho$-orders at least $9$.

Choose $r_c^{11}=\pi$, $s^7=\pi$, and $\rho^{15}=\pi$ compatibly,
and set
\[
 \alpha_7=s^2r_c^{-3},\qquad \alpha_{15}=\rho^{-4}r_c^3.
\]
Then $\alpha_7^{11}=s$, $\alpha_{15}^{11}=\rho$, and
$\alpha_h^h=r_c$.  Rescale the opposite node parameter to obtain
$a_hy_h=\alpha_h$.  This fixes the smoothing units, not just their
valuations.  The central residues of $a_7,a_{15}$ can be chosen
reciprocal.

At the generic point of $N$, the identities
$d\beta=23\beta\eta$ and \cref{eq:harmonic-central-map,eq:harmonic-effective-differential}
give the leading inverse-branch equation
\[
 (\delta z)^{22}=\pi^2z^8(1-z^2)^{22}.
\]
The coefficients of Taylor degrees $1$ through $22$ are divisible
by $23$; the degree-$23$ coefficient reduces to
$-2/(1+z^{23})^2$.  After scaling, the residual equation is a unit
times $v^{23}-v$, so its simple roots label the same $P$ by
$i\in\FF_{23}$.  Its action has leading term
\[
 \sigma_i(z)-z=r_cz^{4/11}(1-z^2)i
       +\text{terms of larger central valuation}.
\]
In the two node coordinates this becomes
\[
 \left.\frac{\sigma_i(a_7)-a_7}{r_c}\right|_C
   =\frac{i}{11}a_7^{-6}(1-a_7^{22}),\qquad
 \left.\frac{\sigma_i(a_{15})-a_{15}}{r_c}\right|_C
   =\frac{i}{11}a_{15}^{-14}(1-a_{15}^{22}).
\]
Apply \cref{lem:harmonic-node-coefficient} to the regular functions
$\sigma_i(a_h)-a_h$.  It gives
\[
 \sigma_i(y_h)-y_h=-\frac{i}{11}y_h^{h+1}+O(y_h^{h+2}).
\]
On the primitive tail $X=y_7^{-11}$, so
$\sigma_i(X)-X=i y_7^{-4}+\cdots$.
On the new tail put $\xi_{\rm tail}=y_{15}^{-11}=X^2/Y$.
Here
\[
 X=y_{15}^{-33}V_\lambda,\qquad
 V_\lambda^5(V_\lambda-1)=\lambda y_{15}^{165},\quad V_\lambda(0)=1,
\]
so $\sigma_i(X)-X=3i y_{15}^{-18}+\cdots$.
Thus the Laurent root labels are $i$ and $3i$ for the same wild
element.  This conclusion uses two vertical valuations of an actual
node function; it is not obtained by exchanging limits in a central
expansion.  Eleventh-root ambiguities multiply labels by squares
in $\FF_{23}^{\times}$ and cannot exchange just one of the designs
in the next calculation.

\begin{computation}[The compatible first jet]
\label{comp:harmonic-first-jet}
Trivialize the smooth quadric formally and use the nine-dimensional
cubic slice specified in \cref{sec:harmonic-reconstruction}.
The first necessary function and logarithmic-differential equations
have exactly two geometric solutions:
\[
 c=c_0\pm\pi C_1\pmod{\pi^2},\qquad
 z_A=1+\pi\pmod{\pi^2},\quad z_B=-1+\pi\pmod{\pi^2},
\]
where $C_1=H^3-5HU^2-7HUV+11HV^2$ and the sign of $\pi$ is fixed
by the first effective-fibre differential normalization
$\eta_1=(z+15z^3+12z^5)\,dz\pmod{z^7}$.
In the common labels above, only the positive solution is compatible
with the two full $M_{23}$-tails.
\end{computation}

Here is the finite test, including why its conclusion is exact.
In the slice the residue, regularity, and differential conditions are
a rank-$18$ linear system on $26$ coefficients (the nine cubic
coefficients, two marked motions, and fifteen numerator coefficients).
On its eight-dimensional solution space, the next function equations
have a degree-two residual ideal.  Its reduced Gr\"obner basis is
triangular with one equation $(t+1)(t+2)=0$ and seven linear
equations.  Thus the two solutions above exhaust the geometric
solutions, not just the $\FF_{23}$-rational ones.

With the fixed translation subgroup $P$ in $S_{23}$ there are two
containing conjugates of $M_{23}$: their number is
$|N_{S_{23}}(P)|/|N_{M_{23}}(P)|=22/11=2$.
They preserve the two cyclic Witt designs $\mathcal B_0,\mathcal B_1$
on $\FF_{23}$, each with $253$ blocks and parameters $S(4,7,23)$.
For reproducibility, $\mathcal B_0$ is given by the weight-seven
words of the binary cyclic Golay code with generator
$x^{11}+x^9+x^7+x^6+x^5+x+1$; $\mathcal B_1$ is its image under
$i\mapsto-i$.  Exact permutation calculations verify their
automorphism groups, their intersection $23{:}11$, and generation
of $A_{23}$ by the two groups.

The first-jet equations give the primitive map
$X^{23}+5X^9+6X^2$ and the two candidate new-tail maps
\[
 Y^3=X^5+\varepsilon,\qquad
 T=2Y(X^6+15\varepsilon X),\qquad \varepsilon=1,-1.
\]
These tails have the independently verified $M_{23}$ groups of
\cref{comp:harmonic-tail-groups}, after scaling over $k$.
For a design $\mathcal B$ form
$I_{\mathcal B}=\sum_{B\in\mathcal B}\prod_{i\in B}X_i$.
If the tail group preserves $\mathcal B$, this is a rational
function of the target, integral on its affine line.  Its pole bound
at infinity is $7/23$ on the primitive tail and $21/23$ on the new
tail, both less than one.  It must therefore be constant.

The exact Laurent calculation excludes $\mathcal B_1$ for the
primitive tail and for $\varepsilon=1$, and excludes
$\mathcal B_0$ for $\varepsilon=-1$.  In the stored scaled series,
the first nonzero remainders are respectively $15t^{1771}$,
$5t^{759}$, and $5t^{759}$; the scaling and root equations are
included in the checker.  The computation solves and checks all
$23$ local roots to precision exceeding these exponents.
A nonzero coefficient disproves constancy.  Conversely, no finite
zero expansion is used to prove constancy: the known $M_{23}$ group
and the exact two-design count identify the remaining design.

Both connected tails in one $G$-cover must act as the same subgroup
on $G/M$.  The negative candidate cannot do so, whereas the positive
one agrees with the primitive tail.  This proves the selection in
\cref{comp:harmonic-first-jet}.  In particular, the common sheet
representation supplies information not contained in the two
abstract group names or in the ramification conductors.

\subsection{From necessary jets to exact incidence}

The remaining local input is finite algebra.  It is convenient to
state exactly what is calculated before invoking the all-orders
theorem.

\begin{computation}[The required marked neighbourhood]
\label{comp:harmonic-required-jets}
For the positive first jet, the necessary map and effective
critical-fibre equations give
\[
 \begin{aligned}
 c&=c_0+\pi C_1+\pi^2C_2+\pi^3(C_3+v_0UV^2)\pmod{\pi^4},\\
 z_A&=1+\pi+9\pi^2\pmod{\pi^3},\\
 z_B&=-1+\pi-9\pi^2\pmod{\pi^3},
 \end{aligned}
\]
where $v_0$ remains an integral parameter and
\[
 \begin{aligned}
 C_2&=6H^2U+5U^3+3H^2V+4U^2V+3UV^2,\\
 C_3&=-H^3-10HU^2+6HUV+2HV^2.
 \end{aligned}
\]
No osculating-incidence equation is imposed in obtaining these jets.
\end{computation}

The calculation uses truncated canonical-ring multiplication, the
two marked quintic-section systems, and the square of the effective
degree-eight ramification factor.  All first-order equation and
marked-point directions are tested: the first variation of the
normalized function is zero.  This ensures that an omitted higher
curve coefficient does not enter an earlier necessary equation.
The first-jet ideal is computed with its degree-two identities;
the next step leaves the $UV^2$ coefficient of $C_2$ free, and the
effective square-fibre equations fix it to $3$.  They leave exactly
the displayed $v_0$ direction at the following order.
The scripts and independent residual-system audit are specified in
\cref{sec:certificates}.

One normalization requires care.  Let $b$ denote the computational
map normalized by $b(0)=1$, and $u$ its actual third branch value.
Locally
\[
 b-u=R^2F,\qquad \overline R=z^8,\qquad
 \overline F=z^7\cdot\text{a unit}.
\]
The vanishing first variation gives $u=1\pmod{\pi^2}$.
The coefficient of $z^{15}$ at order $\pi$ then gives
$R(0)\in\pi^2\mathcal O$.  Since $F(0)\in\pi\mathcal O$,
\begin{equation}
 u-1=-R(0)^2F(0)\in\pi^5\mathcal O.
\label{eq:harmonic-branch-scalar}
\end{equation}
Thus replacing $b-u$ by $b-1$ is legitimate in the earlier calculations
modulo $\pi^5$, and constant scaling never changes $d\log b$.
At all orders, however, $u$ must be eliminated, not set to one.
The fixed Hensel system below does exactly this.  A finite-residue
check finds valuation exactly five for $u-1$ in the reconstructed
lift, so the distinction is genuine.

We have now derived every hypothesis of
\cref{thm:harmonic-local-incidence} for the lift of the constructed
special map.  The integral canonical equations come from
\cref{lem:harmonic-canonical-lattice}; the effective cluster and
differential from \cref{prop:harmonic-effective-differential}; and
the required jets from the common-sheet comparison and
\cref{comp:harmonic-required-jets}.
The remaining theorem identifies any such lift with an independently
verified exact model.  Its proof uses necessary equations and
uniqueness, not a claim that those equations generate the entire
Hurwitz ideal.  Together these results prove
\cref{thm:harmonic-construction-summary}.

\subsection{Exact reconstruction in a specified harmonic neighbourhood}
\label{sec:harmonic-reconstruction}

We isolate the local theorem used to complete the construction.
Its low-order coefficients are explicit hypotheses, derived for our
special map in the preceding subsections.  The proof has two independent
parts: necessary equations have at most one solution, and a small exact
model supplies a solution.  This separation is what permits an
all-orders conclusion without extrapolating from finite precision.

Let $\mathcal O$ be the ring of integers of
$K_{0,23}=\QQ_{23}(\pi)$, where $\pi^2=-23$.
In canonical coordinates $[H:U:V:W]$ put
\[
 q_0=HW-U^2-UV-V^2,\qquad c_0=UW^2-V^3.
\]
The normalization of their common zero curve in characteristic $23$
has, on $W=1$, the parametrization
\[
 (H,U,V)=(z^6+z^4+z^2,z^3,z).
\]
The two marks are $z=1,-1$; the critical cluster is at $z=0$.
In these coordinates the normalized reduced map is
$(1-z^{23})/(1+z^{23})$.  Thus
$t_{\mathrm n}=(1-z)/(1+z)$ relates this description to
\cref{prop:harmonic-position}; the singular point $z=\infty$
has $t_{\mathrm n}=-1$.

Write
\[
 (m_0,\ldots,m_8)
 =(H^3,H^2U,HU^2,U^3,H^2V,HUV,U^2V,HV^2,UV^2)
\]
and set
\begin{align*}
 C_1&=H^3-5HU^2-7HUV+11HV^2,\\
 C_2&=6H^2U+5U^3+3H^2V+4U^2V+3UV^2,\\
 C_3&=-H^3-10HU^2+6HUV+2HV^2.
\end{align*}
The neighbourhood considered below consists of equations and marks
of the form
\begin{equation}\label{eq:harmonic-neighbourhood}
\begin{aligned}
 q&=q_0,\\
 c&=c_0+\pi C_1+\pi^2C_2+\pi^3C_3
       +\pi^3 u_0UV^2+\pi^4\sum_{i=0}^{7}u_{i+1}m_i,\\
 z_A&=1+\pi+9\pi^2+\pi^3u_9,\\
 z_B&=-1+\pi-9\pi^2+\pi^3u_{10},
 \qquad u_i\in\mathcal O.
\end{aligned}
\end{equation}
The remaining coordinates of the marks are determined on the smooth
central chart by $q=c=0$.

\begin{theorem}[Incidence in the specified local neighbourhood]
\label{thm:harmonic-local-incidence}
Suppose a smooth genus-$4$ curve and a degree-$23$ three-point map
admit the canonical equations and ordered totally ramified marks in
\cref{eq:harmonic-neighbourhood}.  Suppose its remaining ramification
consists of eight simple points in the formal neighbourhood of $z=0$,
and that the integral differential
\[
 \eta=\frac{1}{23}\frac{d\beta}{\beta}
 \quad\hbox{has reduction}\quad
 \overline\eta=\frac{2z^8}{z^2-1}\,dz,
 \qquad \operatorname{div}(\beta)=23A-23B.
\]
Then the marked canonical curve is unique up to the coordinate
normalizations just specified.  It is isomorphic over $K_{0,23}$
to the exact model below.  Its two osculating planes meet the curve
in two distinct geometric points.

The same conclusion holds over an unramified complete extension of
$\mathcal O$ with the same congruences.  The exact model has a
degree-$23$ map with geometric monodromy $M_{23}$ and represents the
distinguished member of the seven-cover family.
\end{theorem}

\paragraph{An intrinsic choice of coordinates.}
For $i=0,1,2,3$, the spaces
\[
 H^0(X,\omega_X(-iA-(3-i)B))
\]
are one-dimensional in this neighbourhood, and their generators form
a basis.  This follows from rank-three jet minors and a nonzero
basis determinant on the closed fibre.  Use these generators as
coordinates $x_0,x_1,x_2,x_3$.  The marks become
$[1:0:0:0]$ and $[0:0:0:1]$, with vanishing sequences
$(0,1,2,3)$ and $(3,2,1,0)$.
The quadric consequently has the form
\[
 ax_0x_3+bx_0x_2+dx_1x_2+ex_1^2+fx_1x_3+gx_2^2.
\]
Here $a,b,d,e$ are units.  Substituting $x_i=d_i y_i$ and multiplying
the equation by $\kappa$, with
\[
 d_0=1,\qquad d_1=b/d,\qquad
 \kappa=(ed_1^2)^{-1},\qquad
 d_2=(\kappa b)^{-1},\qquad d_3=(\kappa a)^{-1},
\]
sets the first four displayed coefficients to one.
This normalization is unique and requires no extension of the field.
Reduce the cubic modulo the monic quadric in lexicographic order
$x_0>x_1>x_2>x_3$, and set its $x_0^2x_3$ coefficient to one.
These operations remove the scalar choices in the four generators.

\paragraph{The exact model.}
In these normalized coordinates the quadric is
\begin{equation}\label{eq:small-harmonic-quadric}
 q_*=x_0x_2+x_0x_3+x_1^2+x_1x_2
             -\frac{127}{229}x_1x_3+\frac{12192}{52441}x_2^2.
\end{equation}
The cubic $c_*$ has the following nonzero coefficients; all others
vanish.
\begin{center}
\renewcommand{\arraystretch}{1.18}
\begin{tabular}{@{}ll@{}}
\toprule
Monomial & Coefficient in $c_*$\\
\midrule
$x_0^2x_3$ & $1$\\
$x_0x_1x_3$ & $(69+41\pi)/916$\\
$x_0x_3^2$ & $-(288+3552\pi)/52441$\\
$x_1^3$ & $-(5+\pi)/458$\\
$x_1^2x_2$ & $-(3061+569\pi)/209764$\\
$x_1^2x_3$ & $-(8649+1515\pi)/104882$\\
$x_1x_2^2$ & $-(298479+49803\pi)/48035956$\\
$x_1x_2x_3$ & $-(1041970+181274\pi)/12008989$\\
$x_1x_3^2$ & $(2318419+389759\pi)/48035956$\\
$x_2^3$ & $-(2350008+338328\pi)/2750058481$\\
$x_2^2x_3$ & $-(55642056+9354216\pi)/2750058481$\\
\bottomrule
\end{tabular}
\end{center}

\begin{proof}[Proof of \cref{thm:harmonic-local-incidence}]
We separate local uniqueness, exact existence, and identification of
the cover.  In particular, recognizing coefficients from a local
approximation is not used as a proof of their exact values.

\emph{Local uniqueness.}
The nine cubic directions $m_i$ give a formal slice for canonical
curves with quadric $q_0$: the six infinitesimal orthogonal-coordinate
directions, cubic scaling, the four multiples of $q_0$, and these
nine directions span the twenty-dimensional cubic space.
The determinant is $10$ in $\FF_{23}$ for the certificate's basis.
The formal inverse-function theorem therefore justifies this slice.

Choose bases $S_A,S_B$ of the four-dimensional kernels of the
order-$23$ quintic jet maps.  Fixed unit minors make their
coefficients analytic functions of the eleven parameters $u_i$.
An invertible minor of fifteen cross-product equations determines
a multiplier matrix $M$, with one entry fixed to remove its scalar.
It defines a local ratio $N/D$, normalized to value one at $z=0$,
even off the locus where all global cross-products vanish.
On a genuine cover this is its function $\beta$, up to scalar.
No global map is inferred merely from solving those fifteen rows.

Let $R(z)$ be the degree-eight Weierstrass factor of $\eta$, with
$R\equiv z^8\pmod\pi$.  Reduce $N-D$ and $D$ modulo $R^2$, obtaining
$P,Q$, and put
\[
 E=P-\frac{P(0)}{Q(0)}Q,\qquad E_j=[z^j]E.
\]
Here $Q(0)$ is a unit.  Retaining $R^2$ retains the effective
ramification multiplicities.  For a three-point map $E=0$,
because its eight simple critical points have one common value.
Let $C_j$ denote the central-series coefficient of $z^j$ in the
negative of the first remaining cross-product
$S_{A,0}(MS_B)_1-S_{A,1}(MS_B)_0$, using the fixed bases in the
certificate.  Eleven necessary equations are
\begin{equation}\label{eq:harmonic-hensel-equations}
 F=\left(\frac{E_{15}}{\pi^3},
       \frac{E_9}{\pi^4},\ldots,\frac{E_{14}}{\pi^4},
       \frac{C_7}{\pi^5},\ldots,\frac{C_{10}}{\pi^5}\right)=0.
\end{equation}
These are integral restricted power series on the stated neighbourhood.
The exact residual calculation gives
\[
 \overline F=J\overline u+b,\qquad
 \det J=16\in\FF_{23}^{\times},\qquad
 -J^{-1}b=(0,0,-7,0,-1,-2,0,10,0,-7,-7).
\]
The full matrix, fixed jet bases, and normalizations are retained in
the local certificate and its accompanying derivation.

This calculation is an identity test, not a sample of possible covers.
The parameters first enter at orders $\pi^3,\pi^4$.
The first variation of the normalized ratio vanishes modulo $\pi$;
its variable contribution begins at order $\pi^4$.
After division by $23=-\pi^2$, the coefficients of $R$ first vary
at order $\pi^2$.  Modulo the precisions needed in
\cref{eq:harmonic-hensel-equations}, all expressions have parameter
degree at most two.  Values at zero, the positive and negative
coordinate vectors, and the sums of two distinct coordinate vectors
therefore determine them.  The $78$ evaluations verify integrality
of the divisions and the displayed affine reduction.
The computation also verifies that the two unavailable final digits
of the Weierstrass factor do not affect these divided reductions.

Consequently $F(u)=\widetilde J u+\widetilde b+\pi G(u)$, with
$G$ integral and restricted and $\widetilde J$ invertible over
$\mathcal O$.  The map
$u\mapsto-\widetilde J^{-1}\widetilde b
          -\pi\widetilde J^{-1}G(u)$
is a strict contraction on $\mathcal O^{11}$.
It has a unique fixed point, also after the stated base extensions.
Every cover in the theorem satisfies these necessary equations.
We do not need to assert that the selected equations alone imply
all the other map equations.

\emph{Exact existence.}
The pair $(q_*,c_*)$ was recovered from the local lift in the
intrinsic coordinates above, without using the previously stored
generic models.  The following independent tests are over $K_0$.
The Jacobian-minor ideal is the unit ideal on all four affine
projective charts, so the curve is a smooth genus-$4$ complete
intersection.  Its marked vanishing sequences are as stated.
The order-$23$ quintic jet kernels each have dimension four.
Their common base divisors are exactly $23A$ and $23B$: saturation
excludes other base points, and a section has order exactly $23$
at each mark.

In the degree-ten canonical ring, of dimension $57$, an exact
rank-$15$ multiplier calculation and all six cross-product
identities give a function with divisor $23A-23B$.
Away from its zero and pole, a critical ideal in the chart $x_1=1$
has an eight-dimensional reduced quotient algebra, and $N=uD$
in that algebra for one nonzero scalar $u$.
Reducedness is certified by a squarefree degree-eight characteristic
polynomial of a multiplication operator.
Riemann--Hurwitz gives total ramification $52$; the two marked
points contribute $44$, and the eight distinct critical points
exhaust the remainder.  Thus each is simple, there are no further
critical points in another chart, and the map has the required
three-point passport.

An integral projective change with unit determinant identifies
this exact marked model with the stated local slice modulo
$\pi^{33}$.  Correcting the quadric and cubic by the formal slice
construction preserves the lower-order congruences in
\cref{eq:harmonic-neighbourhood}.  The normalized map is determined
by its divisor; its divided differential and critical cluster are
therefore those used to form the necessary equations above.
The exact map belongs to the required neighbourhood and hence,
by uniqueness, is the local lift to all orders.

\emph{Incidence and identification.}
The two osculating planes on this exact model are $x_3=0$ and
$x_0=0$.  On their common line the identity
\begin{align*}
 c_*(0,x_1,x_2,0)
 &=\left(x_1^2+x_1x_2+\frac{12192}{52441}x_2^2\right)\\
 &\hspace{5mm}\cdot
 \left(-\frac{5+\pi}{458}x_1
             -\frac{771+111\pi}{209764}x_2\right)
\end{align*}
has first factor equal to $q_*|_{x_0=x_3=0}$.
Its discriminant is $3673/52441\ne0$.  This proves the claimed
two reduced points, exactly rather than to finite precision.

Finally, a separate comparison certificate constructs an invertible
projective matrix between this model and the degree-one model of
\cref{sec:hurwitz-scheme}.  It sends the ordered marks to the
ordered marks and verifies equality of the two rational functions
up to a nonzero target scalar by exact reduction in the canonical
ring.  The existing certified branch cycles consequently give
geometric monodromy $M_{23}$ for this map.  This is an exact
isomorphism argument, not an inference from the passport and not
a new numerical continuation calculation.
\end{proof}

\Needspace{8\baselineskip}
\begin{remark}[An explicit descent of the reconstructed map]
The exact model also admits a direct descent, independently of
the exclusion of incidence on the other six covers.
Put $a=-127/229$, $b=12192/52441$, and define
\[
 T(x_0,x_1,x_2,x_3)=(abx_3,bx_2,x_1,x_0/a).
\]
If $\sigma(\pi)=-\pi$, exact substitution gives
\[
 T^2=bI,\qquad q_*(Tx)=bq_*(x),\qquad
 \sigma(c_*)(Tx)\equiv\mu c_*(x)\pmod{q_*},
\]
where $\mu=-(52130760+33951768\pi)/1525141603$ and
$\mu\sigma(\mu)=b^3$.
Thus $R=(b/\mu)T$ satisfies $R\sigma(R)=I$.
Solving $v=R\sigma(v)$ gives a four-dimensional $\QQ$-space
and an invertible change of coordinates in which the curve
equations are rational; the certificate checks these equations
directly.

For the map $f$ normalized to have third branch value $1$, one also
has $f(x)\sigma(f)(Tx)=1$ in the function field.
Consequently $t=(f-1)/(\pi(f+1))$ is fixed by the descent and
defines a map over $\QQ$.  Its branch locus is $t=0$ and
$23t^2+1=0$.  Rescaling by $T_{\mathrm{base}}=23t$ gives the
branch locus $0,\pm\sqrt{-23}$ used in the paper.
This recovers the known rational cover in different equations.
\end{remark}

\begin{remark}[Why uniqueness is enough]
The eleven equations are necessary for a cover in this neighbourhood.
We do not assert that their ideal is the entire map ideal.  Nor do
the displayed low-order incidence cancellations prove an exact
incidence relation on their own: an off-locus audit detects a higher
defect in the constant residual multiplier.
What closes the proof is that the constructed lift and the
independently verified exact model both satisfy the necessary system,
which has at most one solution.  Their identification transfers the
exact incidence calculation to the lift.  The earlier coefficient
checks place the model in the neighbourhood; uniqueness, not the
number of matching digits, supplies equality to all orders.
\end{remark}


\subsection{What the construction explains}

The special map supplies the harmonic configuration, but the
characteristic-zero incidence needs more.  The marked distance fixes
the descent action on the good curve.  Its canonical lattice and
effective ramification divisor fix the closed canonical geometry.
The common Mathieu action excludes the other first jet.  Finally,
unit-Jacobian uniqueness and exact existence identify the entire
lift.  Each step retains information which the preceding, coarser
description would lose.

There are also two distinct uniqueness statements.  Local uniqueness
identifies a lift in this neighbourhood; it does not assert that
all covers with the passport enter this neighbourhood.  Global
uniqueness of the incidence locus is proved on all seven exact
models in \cref{comp:osculating}.  Its intrinsic definition then
gives Galois fixedness by \cref{prop:relative-osculating}.

Accordingly, the construction explains the exceptional incidence for
an explicit harmonic special $M_{23}$-map.  It remains computational
at the tail-group, jet, and exact-model stages.  It is not a
permutations-only existence argument, nor a general theorem that
harmonic reduction forces rationality.  The finite Fano--affine
comparison considered next is a further question, not an unstated
ingredient of this proof.
